\section{はじめに}
検出と命名を分離した粒度可変物体アノテーションでは，Detection Transformer（DETR）\cite{ref_DETR}またはYOLO\cite{ref_YOLO}で検出した領域のクロップに対し，Contrastive Language-Image Pre-training（CLIP）\cite{ref_CLIP}のゼロショット分類で深さ0の単語を決め，WordNet\cite{ref_WordNet}の下位語を候補語彙として同じクロップの分類を繰り返す．
現在の構成には3つの課題がある．
第一に，CLIPの確信度は候補語彙の中で正規化するソフトマックスで計算されるため，候補語彙に正しい語が含まれない場合でも，候補の中で相対的に最もよい語に高い確信度が与えられる．
この結果，深さ0が誤ったまま下位語へ展開された枝において，深い段階の確信度で誤りを検出できない．
第二に，COCO\cite{ref_COCO}の80クラスで学習したYOLOを深さ0に用いるハイブリッド構成は，80クラスの範囲では高い精度を示すが，YOLOの予測を覆す規則はCOCO外の22語の確信度に対する閾値に限られる．
第三に，COCO外の22語は人手で選んだ語であり，YOLOを再学習せずにCOCO外の語彙へ体系的に対応する方法がない．

本稿では，これらの課題に対応する3つの変更として，シグモイドで確信度を計算するSigLIP\cite{ref_SigLIP}への置き換え，YOLOとVLMのスコアの比較による深さ0の決定，およびWordNetにおいてCOCOのカテゴリと同じ深さにある語彙による深さ0の候補を導入し，現在の構成と同一の分割・評価指標で比較する．
